\subsection{The positive primed Witt index}
\label{tpd:main}

Put \(k=(a-a')/2\), and let \(H\) denote a hyperbolic plane of the
appropriate sign. In the step \(q\mapsto q-1\), the unchanged source
is \(V\), the other source changes from \(V'_1\oplus H\) to \(V'_1\),
and the target from \(W=W_1\oplus H\) to \(W_1\). We use the
splitting characters of \S\ref{newtheta:global} throughout this tower.
Write \(P=MN\) for the target parabolic stabilizing the first
isotropic line of \(H\), and normalize \([N]\) to have volume one.
Write \(m(t,h_1)=\operatorname{diag}(t,h_1,\bar t^{-1})\).

\begin{proposition}[The unprimed coefficient when \(p=0\)]
\label{tpd:zero}
Suppose \(p=0\) and \(q>0\). In rational mixed coordinates, let
\(e_0\) be evaluation at zero in the new \(V(\mathbb A_F)\)
variable. For the coherent mixed realizations and the same Haar
measure on the unchanged source, the full, unnormalized constant term is
\[
 (I_\Phi)_P(m(t,h_1)h)=
 \beta(t)|t|_{\mathbb A_E}^{a/2}\,
 I^{W_1}_{e_0(\omega(1,h)\Phi)}(h_1).
\]
Some polynomial archimedean test with an allowable finite factor has
nonzero coefficient. Every unprimed lift with polynomial archimedean
test is orthogonal to the cuspidal spectrum of \(U(W)\).
\end{proposition}
\begin{proof}
In the mixed model the rational summation variables are \((x,y)\),
where \(x\in V(E)\) and \(y\) belongs to a rational position space
for \(V\otimes_E W_1\). The center of \(N\) acts through
\(\psi(bc\langle x,x\rangle_V)\), with \(b\in\mathbb A_F\) and
a fixed nonzero rational factor \(c\in F\) determined by the paired
basis of \(H\). Its integral over \(F\backslash\mathbb A_F\) selects
\(\langle x,x\rangle_V=0\). Definiteness gives \(x=0\).
The other coordinates of \(N\) act on the smaller oscillator through
Heisenberg operators whose translation and modulation parameters are
linear in \(x\), with quadratic central phase zero at \(x=0\).
They therefore act trivially on this value. Both unipotent and source
quotients are compact, and the rational Schwartz sum is absolutely
and uniformly convergent on compact sets of representatives; this
justifies both integrations. The line Levi acts at zero by
\(\beta(t)|t|_{\mathbb A_E}^{a/2}\), the square root of the modulus
of scalar substitution on the \(a\) coordinates times the prescribed
Siegel character. The smaller target uses the same splitting; on the
source the added plane acts on this variable by pullback. This is
also the mixed-model tower formula for a cuspidal source, applied to
the trivial representation of anisotropic \(U(V)\).
\cite[Proposition~3.1, (3.2)--(3.4)]{wuThetaTower}

There are \(q\) rational hyperbolic planes in \(W\). At split finite
places it is split, at inert places it is hyperbolic by the perfect
lattice calculation of Proposition~\ref{threal:allinert}, and at
every real place its Witt index is \(q\). To split off one plane,
multiply the skew form by a nonzero trace-zero element of \(E\) and
apply Hasse--Minkowski to the resulting quadratic norm form on the
underlying \(F\)-space. Complete an isotropic vector to a hyperbolic
pair. Witt cancellation preserves the remaining local indices, so
the argument can be iterated. The bottom \(X\) is definite of
dimension \(d_0=d-2q\).

The unprimed nonvanishing calculation in Theorem~\ref{threal:base}
applies to \((V,X)\): its argument up to the evaluation of
\(Z((d_0-a)/2)\) uses only these two spaces. Indeed
\(a'-2q>1\), \(d_0\) is even and the infinite angular exponent still
exceeds \(d_0\). At inert places \(X\) is hyperbolic and admits a
self-dual lattice, and at split places one may take dual bases.
These lattices agree almost everywhere with the lattice obtained by
a rational orthogonal decomposition. The denominator arguments
are at least \(a'-2q+1>1\) and the numerator integers are
\(d_0/2-i\) for \(0\leq i<a\); the same functional equation and
nonvanishing at one apply. At infinity the unprimed invariant
determinant used there has exponent \((m_\alpha-m_\beta-d_0)/2>0\)
and there are \(d_0\geq a\) columns. This gives a nonzero bottom
period.

Evaluation in each new variable is onto the smaller Schwartz
space by tensoring with a test of value one at zero. These extensions
are restricted tensors, since away from finitely many places the
chosen rational hyperbolic basis and its inverse are integral. The
constant-term formula gives nonvanishing successively in the tower.
With the finite test fixed, the resulting nonzero coefficient at
a fixed Levi element is a continuous functional of the infinite
Schwartz test, by compactness of the integrating quotients.
Polynomial times Gaussian tests in the original coordinates are
dense in the Schwartz topology, giving the asserted polynomial test.

For orthogonality we keep the unconjugated convention of Wu, in which
the two directions use inverse additive characters. Set
\(Y=V\), and let \(\chi_1,\chi_2\) be the splitting characters on
the target \(U(X)\) and the source \(U(Y)\), respectively, in
Wu's labeling (the twists below mean composition with determinant).
The bottom lift \(\sigma=\theta^X_{Y,\psi}(\mathbf1)\) is nonzero
and cuspidal since \(X\) is anisotropic. The input is itself an
irreducible cuspidal representation. Wu's Theorem~5.5 gives
irreducibility of \(\sigma\) and Theorem~5.4(1), with the two
additive characters in their indicated order, gives
\[
 \theta^Y_{X,\psi^{-1}}(\chi_1^{-1}\sigma)=\chi_2\mathbf1.
\]
The input and output in this last equality are again nonzero,
irreducible and cuspidal. Applying Theorem~5.1(2) with
\(\pi=\chi_1^{-1}\sigma\) therefore shows that
\(\theta^{X\oplus H^q}_{Y,\psi}
 (\chi_2^{-1}(\chi_2\mathbf1))\) pairs to zero with every cusp
form for \(q>0\). This is the asserted unprimed lift.
\cite[Theorems~5.1, 5.4 and 5.5]{wuThetaTower}
\end{proof}

\begin{proposition}[Cuspidal orthogonality of the supported primed lift]
\label{tpd:orth}
For every \(q>0\), including \(p>0\), every supported primed
theta integral pairs to zero with all cuspidal forms on \(U(W)\).
\end{proposition}
\begin{proof}
First take a compact-finite automorphic vector \(u\) in a single
irreducible cuspidal constituent on \(J=U(W)\). Such vectors span
a dense subspace of the cuspidal spectrum. Set
\[
 B(g)=\int_{[J]}\Theta_{\Phi'}(g,h)\overline{u(h)}\,dh,
        \qquad g\in U(V')(\mathbb A_F).
\]
Use Siegel representatives for \(h\) in this integral. Rapid decay of
\(u\) there dominates every fixed positive power of \(H(h)\), by
the rapid decrease theorem for cusp forms. Proposition~\ref{threal:decay}
then permits Fubini, and \(B\) is a rapidly decreasing, square-integrable
reverse theta lift. More explicitly its proof gives a bound
\(C_{L,u}\prod_i T_i^{-L}\) on each source Siegel set, with the
same bound when the right compact set of representatives is enlarged.
The finite supported condition is preserved by translations in either
group, so the estimate also applies to every fixed right translate.

Consider any rational maximal parabolic of \(U(V')\), stabilizing a
totally isotropic \(j\)-plane. In the reverse direction the input
\(\bar u\) is cuspidal. Proposition~3.1 of Wu applies to this
\(j\)-plane, with the roles of the Hermitian and skew-Hermitian
spaces exchanged (equivalently rescale the two forms by reciprocal
trace-zero scalars). Its unnormalized constant term at
\(m(t,h_0)g_0\), where \(t\in\mathrm{GL}_j(\mathbb A_E)\), is a
function of \(h_0,g_0\) times
\(\xi(\det t)|\det t|_{\mathbb A_E}^{d/2}\) for a unitary
splitting character \(\xi\). Take \(t\) to be the scalar \(R I_j\)
at each real place and the identity at finite places and let
\(R\to\infty\). On this positive chamber ray the integral over the
unipotent quotient tends to zero: its representatives can be taken
in a compact set which stays bounded after conjugation by
\(m(t,1)^{-1}\) on the left and \(m(t,1)\) on the right, and the
preceding bound applies. The modulus of the displayed character
grows as \(R^{[F:\mathbb Q]jd}\), forcing its coefficient to
vanish for every \(h_0,g_0\). Thus \(B\) is cuspidal.
\cite[Proposition~3.1]{wuThetaTower}

The group \(U(V')\) has a proper rational parabolic. Indeed the
given real Witt index is \(q>0\); at inert places the perfect
Hermitian lattice has Witt index \(\lfloor a'/2\rfloor\), and at
the remaining finite places it is split. The same norm-form
argument as above supplies a rational isotropic line.
The orthogonal projection in \(L^2([U(V')])\) of the constant
function to the closed cuspidal subspace is right invariant, hence
constant. Its constant term at the line parabolic must vanish, so
this projection is zero. It follows that \(\int B(g)\,dg=0\).
Fubini identifies this with the asserted pairing.
\end{proof}

\begin{lemma}[Removal of the primed zero orbit]\label{tpd:primezero}
Suppose \(p=0<q\). In taking the constant term of the supported
primed integral along \(P\), one can choose a test giving the same
automorphic function for which the \(x=0\) summand vanishes
identically, also after any target translation.
\end{lemma}
\begin{proof}
At one real place average the upper test over \(K_+=U(q)\), the
positive compact block of \(U(V')\). The unchanged source integral
follows by right invariance of Haar measure, compact averaging and
Proposition~\ref{threal:decay}. The full-rank condition at \(v_1\)
is preserved. Mixed evaluation at \(x=0\) is \(U(V')\)-equivariant
to the smaller oscillator for \((V',W_1)\): on the extra variable
the group acts by pullback, using the coherent splitting on decreasing
the partner by a hyperbolic plane. Thus its value on the averaged
test belongs to the \(K_+\)-invariants of that oscillator.

Write \((r_1,t_1)=(d-q-1,q-1)\) for the signature of \(W_1\).
As in the entry calculation of Proposition~\ref{threal:allreal},
the \(K_+\)-vacuum character in the primed oscillator is
\(\det^{b}\), where \(b=(m_\alpha-m_\beta+r_1-t_1)/2>0\).
There are \(r_1\) standard and \(t_1\) dual column variables
for this block. By the Cauchy decomposition the compact types can
only come from
\(\det^b\otimes S_\lambda(\mathbb C^q)
                    \otimes S_\mu(\mathbb C^q)^*\),
with \(\ell(\mu)\leq t_1=q-1\); variables from other rows are
trivial on \(K_+\). An invariant would equate the highest weights
\(\mu=\lambda+(b,\ldots,b)\), impossible at the last entry.
Hence the compact average is zero on the polynomial module, and
by density and continuity on the Schwartz space. Target translations
commute with compact source averaging. The smaller evaluation still
vanishes after applying any source element by equivariance.

On a compact truncation of the integrating source, average the
center of \(N\) first. Additive orthogonality leaves the sum over
the rational null cone, and the zero summand just treated vanishes
pointwise. Taking an exhaustion of the source quotient gives the
constant term as the limit of these nonzero-orbit expressions;
the original kernels are dominated by Proposition~\ref{threal:decay}
uniformly over compact representatives for \([N]\). This formulation
requires no separate integrability assertion for an unrestricted
smaller zero-orbit period.
\end{proof}

\begin{corollary}[Detection for a definite unprimed source]\label{tpd:detect}
If \(p=0,q>0\), two matched polynomial tests give proportional
automorphic functions with scalar \(c\) as soon as their full line
constant terms are proportional with that scalar at every target
translate. For the unprimed form this constant term is the smaller
period of Proposition~\ref{tpd:zero}; for the primed form it consists
only of the nonzero orbit of Lemma~\ref{tpd:primezero}.
\end{corollary}
\begin{proof}
For any \(q\) this follows from Proposition~\ref{threal:detectline}.
We give a direct argument in rational Witt index one.
Suppose \(q=1\).
The rational Witt index of \(W\) is one by its signatures and the
isotropy proved above. Every proper rational parabolic stabilizes a
line and these lines are rationally conjugate by Witt's theorem. The
difference \(D\) of the forms thus has every proper constant term zero.
Here one can see square integrability directly. The seminorm bounds
in Proposition~\ref{threal:decay} give a single moderate growth
exponent for all right archimedean derivatives of \(D\). Indeed
choose the source decay order in that proof once, large enough for
source integration. Applying a right derivative replaces the infinite
test by its derived oscillator action. The orders of the seminorms
used to bound the translated test and the powers of the target
height in those bounds depend on this fixed decay order and the
dimensions. Only the seminorm values of the initial test change.

On a rank-one Siegel set write the target representative as \(tc\)
with positive chamber coordinate \(R\ge1\) and bounded \(c\).
At fixed finite level the compact unipotent quotient in its constant
term reduces to a compact real nilmanifold. Indeed strong
approximation on the unipotent group follows successively for its
additive root spaces, and the finite open stabilizer can be taken
common to the bounded finite components. On this nilmanifold the
function \(n\mapsto D(ntc)\) has mean zero. Let \(\Delta\) be
the sum of negative squares of a left-invariant real basis; Haar
volume is invariant and its kernel consists of constants. Elliptic
Sobolev embedding and the spectral gap on the mean-zero subspace
bound a point value by \(C_j\|\Delta^jD(\cdot\,tc)\|_{L^2}\)
for every sufficiently large \(j\). Each basis derivative becomes
a right derivative on \(D\) with Lie algebra argument
\(\operatorname{Ad}((tc)^{-1})X=O(R^{-1})\), on choosing the
coordinate with smallest positive root \(R\). The uniform growth
bound on the compact set of unipotent representatives now gives
\(D(tc)=O_j(R^{B-2j})\) for a fixed \(B\). Thus \(D\) is rapidly
decreasing and belongs to the cuspidal subspace of \(L^2\).
Propositions~\ref{tpd:zero} and~\ref{tpd:orth} make it orthogonal
to the cuspidal spectrum, hence to itself.
\end{proof}

\begin{lemma}[The Mellin exponent and the good Euler factors]
\label{tpd:mellin}
Choose paired bases \(x,y\) and \(x',y'\) of the source and target
hyperbolic planes, respectively. In the rational mixed polarization
for simultaneous removal of these two planes, take the following
isotropic space to act by modulations:
\[
 S=Ex\otimes_E W\ \oplus\ V'_1\otimes_E Ex'.
\]
Write \(e'\) for evaluation at zero in its dual position coordinates,
and \(\rho\) for the Heisenberg action. With \(K\) the product of the
specified source compact stabilizers, the integral
\[
 T\Phi'(z_1)=\int_K\int_{\mathbb A_E^\times}
   \big[e'\rho(\bar l\, y\otimes x')\omega'(k_0,1)\Phi'\big](z_1)
           \,\nu(l)|l|_E^{k+1}\,d^\times l\,dk_0
\]
is absolutely convergent and defines a smaller Schwartz function.
It preserves the required full-rank support at \(v_1\).
Its target line character is \(\beta(t)|t|_E^{a/2}\).
With local unit groups of multiplicative mass one almost everywhere,
the scalar from the unmodified unramified factors outside any
sufficiently large finite set is the corresponding partial Euler
product for \(L_E(\nu,k+1)\).
\end{lemma}
\begin{proof}
The chosen subspace has dimension \(d+a'-2\) over \(E\);
its symplectic perpendicular modulo itself is the smaller tensor
space \(V'_1\otimes_EW_1\), by the two orthogonal decompositions.
We can choose its dual coordinate space to contain the translation
line \(E(y\otimes x')\), paired nondegenerately with
\(E(x\otimes y')\). Consequently the bracket in the integral is the
restriction of a mixed-model Schwartz function to a linear subspace
with coordinates \((\bar l,z_1)\), allowing a fixed sign or unit
change in the first coordinate. A rational choice of these
coordinates is integral with integral inverse almost everywhere,
where the vacuum is a nonzero multiple of the product lattice
indicator and the compact source fixes it. Choosing each original
basic vector and its image as the smaller basic vector removes the
same multiple on both sides. For a fixed finite test the compact
finite orbit is finite, by smoothness; at infinity the compact
orbit is bounded in every Schwartz seminorm. Finite-place restrictions
are finite sums of products of locally constant compactly supported
functions of \(l\) and \(z_1\) off the basic places. At infinity
all smaller Schwartz seminorms and derivatives in \(z_1\) are
dominated by arbitrarily rapidly decreasing majorants in \(l\).
Since \(k+1>1\), the integrals at zero converge. At a basic place
\(w\) of \(E\) the integral is the geometric series
\(\sum_{j\ge0}\nu_w(\pi_w)^j(Nw)^{-j(k+1)}\) times the smaller
basic vector. The absolute products converge since the Dedekind
zeta series converges at \(k+1\). This proves the analytic and Euler
factor assertions, including differentiation under the integrals.

At \(v_1\), choose the two conjugate idempotent factors of the
split tensor space as position and momentum in the specified primed
matrix model. Order the rows by \((x,V'_1,y)\) and the columns
by \((x',W_1,y')\) on the position side. The complementary momentum
basis pairs \(x\) with \(y\) and \(x'\) with \(y'\). In passing to
the above mixed evaluation, the position vectors belonging to \(S\)
are integrated out by partial Fourier transformation. They are the
first row and the first column of the middle rows. The momentum
vectors belonging to \(S\) require zero in the last row and the last
column of the middle rows. Translation by \(\bar l\,y\otimes x'\)
puts a nonzero entry in the first position of the last row; its other
idempotent component contributes only a Fourier phase. Thus, for
middle position matrix \(Z\), the original test is evaluated solely
on matrices of the form
\[
   \begin{pmatrix} *&*&*\\ *&Z&0\\ b&0&0\end{pmatrix},\qquad
                                      b\in F_{v_1}^{\times}.
\]
Changing the chosen idempotent uses the corresponding conjugate
bases and the same pattern. Clearing the middle entries of the
first column by the last row bounds the full rank by
\(2+\operatorname{rank}Z\). The upper supported test, also after
the compact source action by invertible row operations, vanishes
on all these matrices if \(\operatorname{rank}Z<a'-2\).
Every partial Fourier integral then vanishes at \(Z\), proving
the smaller support assertion.

At infinity an adapted Cartan involution preserves the orthogonal
decomposition into the tensor blocks formed from the two planes
and their complements. In particular the smaller symplectic tensor
is a complex-orthogonal summand for the compatible positive
structure. The Fock vacuum factors over this summand and its
complement; the polynomial module is spanned by finite products of
their creation monomials. Compact source translates of a polynomial
test remain in the finite span of fixed degrees. The evaluation and
translation in the integral take place only in the complement.
The smaller real output is therefore polynomial for polynomial
input, in the coordinates transported by the smaller Cartan
structure. The moments in the complement converge by the same
Gaussian bounds.

Furthermore fixing a nonzero common pair of real matching maps fixes
the infinite component of \(\nu\). Given two allowed systems for
these same maps, their primed source actions on a compact unitary
block at one real place differ by the determinant character
\(\det^{(m_\nu-m_{\widetilde\nu})/2}\), by the doubled splitting
calculation in Proposition~\ref{threal:allreal}. Both actions leave
the same quotient map unchanged. Apply that map to a polynomial
with nonzero image and translate by this compact block. The
determinant character is identically one since the block determinant
ranges over the circle, so \(m_\nu=m_{\widetilde\nu}\) there.
The restriction of each character to \(\mathbb A_F^\times\) fixes
its radial part at infinity. Varying the place proves the claim.
In particular the real Mellin operator itself is fixed across
the allowed characters and actions for fixed real maps, after the
same coordinates and Haar measures are chosen.

The source line Levi \(m(l,1)\) has determinant modulus
\(|l|_E^d\) on \(S\), since it scales the first summand by \(l\).
The corresponding evaluation character is \(\nu(l)|l|_E^{d/2}\).
Indeed first split the source plane, with modulation space
\(Ex\otimes W\) and its prescribed Siegel character \(\nu\),
then split the target plane in \(V'_1\otimes W\). On the source
unipotent radical the line Levi acts on \(a'-2\) coordinates over
\(E\) with determinant modulus \(|l|_E^{a'-2}\), and on the central
\(F\)-coordinate by \(l\bar l\), adding the factor \(|l|_E\).
Its inverse modular character in Iwasawa integration is therefore
\(|l|_E^{1-a'}\); together these give \(d/2+1-a'=k+1\).
The target line has determinant modulus \(|t|_E^{a'-2}\) on \(S\):
its action on the full \(W\) in the first summand has modulus one.
At evaluation its scalar is \(\alpha(t)|t|_E^{(a'-2)/2}\), from
the residual target Siegel action (keep the target splitting
character on reducing the source by a plane). Moving this target
element past the translation replaces \(l\) by \(l/\bar t\).
In the Mellin integral substitute \(l=u\bar t\). Since
\(\nu(\bar t)=\nu(t)^{-1}\), the resulting phase is
\(\alpha(t)/\nu(t)=\beta(t)\) and the exponent is
\((a'-2)/2+k+1=a/2\), as claimed. This computation explains both
the choice of the translation line and the shift by one in this
branch.
\src{src/Development/AbelianVarietyTheory/Definitions/Divisors/CTSmallCone.lean}{57}{The two-plane evaluation, cone and translation};
\src{src/Development/WeilBound/ThetaLifts/SmallCones/ConeLawAB.lean}{115}{The primed two-Levi character}.
\end{proof}

The two fractional-ideal exponents in this branch can also be fixed.
By adaptedness the line \(Ex\) has lattice \(\mathfrak b x\) and
the line \(Ex'\) has lattice \(\mathfrak c x'\). The source lattice
is Hermitian self-dual and the two lines and the orthogonal smaller
lattice split integrally. Pairing with \(\mathfrak b x\) therefore
shows that the lattice of \(Ey\) is \(\bar{\mathfrak b}^{-1}y\).
It follows by taking direct summands in the tensor lattice that the
finite translations \(\bar l\,y\otimes x'\) belonging to it have
\[
                l\in\mathfrak a\widehat{\mathcal O}_E,qquad
                \mathfrak a=\mathfrak b^{-1}\bar{\mathfrak c}.
\]
Both ideals \(\mathfrak b,\mathfrak c\) are prime to the conductor
of \(\nu\). At the chosen half of a ramified split place, this
follows because both paired vectors occur in the specified integral
basis. At the conjugate half, apply duality to the second vector in
that basis; the target's trace codifferent is a unit at the
designated totally split rational primes. At unramified places
the value of \(\nu\) on a uniformizer is independent of its unit
multiple, and \(\nu(\bar l)=\nu(l)^{-1}\). Consequently
\(\nu(\mathfrak a)=\nu(\mathfrak b)^{-1}
\nu(\mathfrak c)^{-1}\). Integrating the line-lattice indicator
outside the conductor contributes this character value, its absolute
norm to the power \(-(k+1)\), and the elementary Euler factors.
The norm power depends only on the fixed lattice and ranks.
Thus the lattice calculation gives \(e_1=e_2=-1\).
\src{src/Development/WeilBound/ThetaLifts/LocalTheory/ConeShifts/IdealId0.lean}{49}{The character on the translated lattice line}.

\begin{proposition}[Unfolding the supported nonzero orbit]
\label{tpd:unfold}
Suppose \(q>0\). In the rational mixed models the
\(\beta|\cdot|_E^{a/2}\) component of the line constant term of
any supported primed test equals
\[
 C_1\,\beta(t)|t|_E^{a/2}
 I'_1\bigl(T(\omega'(1,h)\Phi')\bigr)(h_1).
\]
Here the line Levi element is \(m(t,h_1)\). The difference from
the entire constant term transforms by \(\alpha|\cdot|_E^{a'/2}\).
It vanishes identically if \(p=0\).
For self-dual Fourier changes \(C_1>0\) depends only on the quotient
and Iwasawa measures and the fixed spaces. It is independent of the
characters and the matching maps.
\end{proposition}
\begin{proof}
We first compute at the target identity for a supported test
whose zero evaluation in the target-first model vanishes identically.
Write this condition as \(e_{W,0}\Phi'=0\). When \(p=0\) the real
average of Lemma~\ref{tpd:primezero} supplies such a test and leaves
both the upper integral and \(T\Phi'\) unchanged, the latter by
right invariance of Haar measure on \(K\).
Put \(Y=V'\), \(Y_1=V'_1\). At a fixed
source element \(g\), integration over the target unipotent center
in the target-first model leaves only null vectors in \(Y(E)\).
Their nonzero terms group by their rational lines. Fix one such
line \(Ex\); its nonzero vectors are \(u x\), \(u\in E^\times\).
Having chosen its companion \(y\), polarize the remaining tensor
\(Y\otimes_EW_1\) using \(Ex\otimes_E W_1\) as modulation space
and a fixed position space for \(Y_1\otimes_EW_1\).
For this vector \(ux\), the remaining target radical acts through
the Heisenberg modulation paired with \(ux\otimes w\),
\(w\in W_1(\mathbb A_F)\). Integrating it over the additive quotient
sets the corresponding rational \(W_1\) position to zero, since
\(u\ne0\). This is the character orthogonality in
\cite[(3.15)--(3.20)]{wuThetaTower}. The line contributes
\(\sum_{u\in E^\times}\sum_{z_1} A_g(u,z_1)\), where \(z_1\)
runs over the rational smaller positions and \(A_g\) is the
restriction of the resulting Schwartz test to the indicated positions.
The source element acts on the test before these operations.

There are now two natural polarizations of the removed coordinates.
The modulation space of this last evaluation is
\[
 S_A=Y\otimes Ex'\ \oplus\ Ex\otimes W_1 .
\]
Its common subspace with \(S\) of Lemma~\ref{tpd:mellin} has
codimension one over \(E\) in both. The two remaining dual directions
are \(y\otimes x'\) and \(x\otimes y'\). Thus replacing \(S_A\)
by \(S\) is a partial Fourier transform in a single \(E\)-coordinate
and leaves \(z_1\) unchanged. Explicitly, on restricting all common
evaluation coordinates to zero, its Fourier transform is
\[
 B_g(l,z_1)=[e'\rho(\bar l\,y\otimes x')\omega'(g,1)\Phi'](z_1).
\]
The trace pairing on the two lines is a nondegenerate rational
pairing; a fixed rational change of variable or sign permutes
\(E^\times\) and has adelic modulus one. With self-dual measures
Poisson summation therefore identifies the sums over \(E\) of
\(A_g\) and \(B_g\), separately at each \(z_1\).
Both missing zero terms vanish. For \(A_g(0,z_1)\) this follows
from \(e_{W,0}\Phi'=0\), since this evaluation intertwines the
source action before the second polarization and evaluation.
For \(B_g(0,z_1)\) use the
matrix in Lemma~\ref{tpd:mellin} with \(b=0\): its last row is
identically zero at \(v_1\). This remains so for the test translated
by \(g\), since source translations preserve full-rank support.
Consequently the sums over \(E^\times\) agree. At fixed \(g\)
the original grouping over lines is legitimate by absolute
Schwartz summation and compact integration in the target radical.

To unfold these line sums set \(D_x=\operatorname{Stab}_{U(Y)}(x)\).
It is the semidirect product of its line unipotent radical \(N_Y\)
and \(U(Y_1)\). The source group is transitive on its rational
nonzero null vectors by Witt extension. In the new expression put
\(f(g)=\sum_{z_1} B_g(1,z_1)\). This function is invariant under
\(N_Y(\mathbb A_F)\): an element of this group changes \(y\)
by a vector in \(Ex+Y_1\), so the change in the translation lies
in \(S\). Its Heisenberg commutator is trivial because \(Ex'\)
is isotropic; evaluation itself is invariant under the unipotent.
For the last assertion split the source plane first: the unipotent
acts on the smaller oscillator through Heisenberg operators with
parameters linear in its extra position, and its center by a
quadratic phase, all trivial at zero. The following evaluation on
the smaller target plane preserves this invariance.
The \(U(Y_1)(F)\) action preserves the rational summation.
For the source line Levi \(m_Y(l,g_1)\), moving it past
\(\rho(y\otimes x')\) replaces \(y\) by \(\bar l\,y\).
The simultaneous evaluation therefore gives
\[
 f(m_Y(l,g_1)k_0)=\nu(l)|l|_E^{d/2}\,
 \Theta_1'\bigl(g_1,1;
    e'\rho(\bar l\,y\otimes x')\omega'(k_0,1)\Phi'\bigr).
\]
For rational \(l\) the scalar is one. Hence summing over rational
lines and their nonzero scalar parameters gives exactly the
periodization of \(f\) over \(D_x(F)\backslash U(Y)(F)\).

We check absolute integrability of this periodization at the level
of the signed smaller theta sum. Integrate first on
\(D_x(F)\backslash U(Y)(\mathbb A_F)\). Iwasawa decomposition for
the source line and the chosen maximal compact gives, up to fixed
positive normalization, integration over \([N_Y]\),
\([U(Y_1)]\), \(\mathbb A_E^\times\) and \(K\), with density
\(|l|_E^{1-a'}\). The finite compact is in good position: it is
hyperspecial for the self-dual source lattice at the inert places,
and is an integral general linear compact at the split places.
The quotient \([N_Y]\) is compact with fixed volume. The absolute
integral remaining is bounded by the integral of the absolute
smaller theta kernel against \(|l|_E^{k+1}\). At the distinguished
finite place each smaller test in this integral has full-rank
support by Lemma~\ref{tpd:mellin}. More explicitly the compact
locally constant test in \((l,z_1)\) has only finitely many slices
in the first variable at each exceptional finite place. At \(v_1\)
all slices occurring for invertible \(l\) have the required support.
At infinity all seminorms needed in Proposition~\ref{threal:decay}
for the smaller source are rapidly decreasing in \(l\) and uniformly
bounded near zero and in \(k_0\). The smaller ranks satisfy its
hypotheses, since \(q\) has decreased by one and \(a'\) by two.
The absolute Euler products outside the exceptional primes are
bounded by the corresponding factors of \(\zeta_E(k+1)\), as
in Lemma~\ref{tpd:mellin}. These bounds prove the required
integrability, justify unfolding the periodization, and allow the
smaller period to be interchanged with the Mellin and compact
integrals. This gives \(C_1 I'_1(T\Phi')(1)\).

Write \(A(t)=\alpha(t)|t|_E^{a'/2}\),
\(B(t)=\beta(t)|t|_E^{a/2}\), and let \(J(\Phi')\) be the entire
constant term at the identity. For arbitrary \(p\) and arbitrary
supported \(\Psi\), apply the computation just proved to
\[
   \Phi'_t=\omega'(1,m(t,1))\Psi-A(t)\Psi .
\]
This test satisfies \(e_{W,0}\Phi'_t=0\): in the target-first
model scalar substitution on the \(a'\) coordinates at zero gives
exactly \(A(t)\). Both summands retain the required split support.
Writing \(J_B(\Psi)=C_1 I'_1(T\Psi)(1)\), the result is
\[
 J(\omega'(1,m(t,1))\Psi)-A(t)J(\Psi)
                   =(B(t)-A(t))J_B(\Psi),
\]
using the Mellin covariance in Lemma~\ref{tpd:mellin}. Thus
\(J-J_B\) transforms by \(A\), whereas \(J_B\) transforms by \(B\).
These components are unique since \(a>a'\) and the splitting
characters are unitary. If \(p=0\), averaging any test first shows
directly that \(J=J_B\). Taking \(\Psi=\omega'(1,h)\Phi'\) and
using intertwining for the smaller target factor gives the asserted
formula at \(m(t,h_1)h\). All measure constants occurred in the fixed
Haar decompositions or unipotent quotients.
\end{proof}

\begin{lemma}[A supported representative for the other coefficient]
\label{tpd:so}
Let \(q>0\). For each supported primed test \(\Psi\), one may
replace \(\Psi\) by a supported \(\widetilde\Psi\) having the
same upper period, so that its one-plane zero evaluation
\(e_{W,0}\widetilde\Psi\) is full-rank supported for the unchanged
source \(V'\) and the target \(W_1\). In the coherent one-plane
model the \(\alpha|\cdot|_E^{a'/2}\) component at \(m(t,h_1)\)
is
\[
   \alpha(t)|t|_E^{a'/2}
       I_{V',W_1}(e_{W,0}\widetilde\Psi)(h_1).
\]
For a fixed further right target translate apply the construction
to the already translated upper test.
\end{lemma}
\begin{proof}
Put \(m=a'\) and work at \(v_1\) in the primed row coordinates.
Choose dual target bases with paired first and last columns. In
this model \(e_{W,0}\phi(Z)=\int\phi(u,Z,0)\,du\), with the
corresponding conjugate choice if the position idempotent is
interchanged. In fact one of the removed columns acts by translations
and is integrated, while the paired modulation column must vanish.
An upper test of full row rank therefore maps to a compactly
supported function on the open region \(D=\{\operatorname{rank}
Z\ge m-1\}\). The map is onto that space. At a given \(Z_0\in D\)
choose \(u_0\) completing its columns to rank \(m\), and small
compact open balls in \(u,Z\) on which rank stays full. A unit-mass
indicator in the first ball and a small last-column indicator of
value one at zero lift every test supported on the second ball.
A finite disjoint refinement covers the support of any test in \(D\).

Modulo compactly supported functions on the open full-rank part
of \(D\), these lower tests restrict to all compactly supported
functions on its relatively closed rank \(m-1\) stratum. Extension
follows by compact open local charts and disjoint refinement.
The stratum's source coinvariants are zero by the rank stabilizer
calculation of Proposition~\ref{split:kernel}, applied with
\((d-2)/2=m+k-1>m\). This calculation concerns the stratum itself
and applies with the same unitary phase and the lowered modulus.
Thus, up to a full-rank lower test, \(e_{W,0}\phi\) is a finite sum
of lower source-action differences of tests in \(D\). Lift each
such test through the surjection above, and subtract the
corresponding upper differences from \(\phi\). Equivariance of
evaluation for the unchanged source proves the required support.
Use individual supported local factors and extend by linearity;
at the other places all factors stay fixed. These upper differences
integrate to zero by right translation on the source, with absolute
convergence from Proposition~\ref{threal:decay}.

In the expansion of the upper constant term the zero-vector
contribution of this representative now converges by the
full-rank estimate applied to \((V',W_1)\). Its torus character is
\(\alpha(t)|t|_E^{m/2}\); at \(t=1\) evaluation gives precisely
the smaller displayed period, using the same source measure.
For each fixed source translate the \(A_g(0,z_1)\) of
Proposition~\ref{tpd:unfold} vanishes even though
\(e_{W,0}\widetilde\Psi\) need not be zero. In fact apply the
further source-plane evaluation \(Ex\otimes W_1\) at the split
place to this lower full-rank test. All entries in one row must be
zero and the opposite row is integrated, so every value in the
integrand vanishes. This remains true for any rational null line,
by changing row basis, and for translates by the unchanged source
and smaller target, which preserve matrix rank. The other Poisson
zero term still vanishes on the upper full-rank test. The entire
nonzero orbit therefore has the \(\beta\) law by the same absolutely
convergent unfolding. Uniqueness of the two characters proves the
assertion, including independence of the resulting component from
the choice of the lift of the source differences.
\end{proof}

At a designated split place the upper to lower comparison space is
at most one-dimensional by Proposition~\ref{threal:localcone}.
Its unprimed zero-evaluation map is nonzero, since evaluation is
onto and the lower quotient is nonzero. We next compute the scalar
on one pair of tests with nonzero smaller image. This determines
the comparison on every test at this place.

\begin{proposition}[The ramified primed comparison on all tests]
\label{tpd:ram}
At a designated ramified split place let \(K_0=F_v\) have residue
size \(Q\), let \(\chi=\nu_w\) on its chosen factor of \(E\),
and let \(c>0\) be the conductor exponent of \(\chi\).
Use the prescribed bases and row polarizations, \(\operatorname{vol}
(\mathcal O)=1\) for the additive character of conductor \(\mathcal
O\), and give \(\mathcal O^\times\) and the maximal source compact
masses one. Use the integral-vacuum-normalized evaluations. On
every upper pair matched at this place, the ratio of the lower
primed quotient on \(T_v\phi'\) to the lower unprimed quotient on
\(e_{0,v}\phi\), as a scalar identity of maps, is
\[
   \epsilon_v(-k)=
      Q^{ck}\chi(\pi)^c
       \sum_{r\in(\mathcal O/\pi^c)^\times}
                         \chi(r)^{-1}\psi_v(r/\pi^c).
\]
\end{proposition}
\begin{proof}
Put \(m=a'\) and \(D_\chi(u)=\mathbf1_{\mathcal O^\times}(u)
\chi(u)^{-1}\), extended by zero at \(u=0\).
For an \(n\)-by-\(d\) matrix choose a block of \(n\) pivot columns.
Let its test be \(D_\chi\) of the pivot determinant times the
indicators that the pivot entries belong to \(\mathcal O\) and
all other entries to \(\pi^c\mathcal O\). Take the first \(a\)
columns for \(\phi\), the last \(m\) for \(\phi'\), and the same
rule for \(\phi_1,\phi'_1\) after deleting the first and last of
the \(m\) target columns. The source compact fixes each test:
integral invertible row operations preserve its entry conditions,
and the source factor cancels the inverse determinant character.

The upper tests are matched, as are the two lower tests. To check
this directly, restrict the orbital functions
\eqref{newtheta:big}--\eqref{newtheta:small} to
\(h=\left(\begin{smallmatrix}A&B\\ C&D\end{smallmatrix}\right)
\in\mathrm{GL}_d(\mathcal O)\).
A contributing row transformation must belong to the integral
general linear group: the original row frame is primitive and
the pivot of the transformed test is invertible modulo \(\pi\).
On both sides support requires \(C=0\bmod\pi^c\), in which case
the diagonal blocks are invertible. The determinant character of
the row transformation cancels throughout that integral compact.
The two values are \(\beta_w(\det h)\chi(\det A)\) and
\(\alpha_w(\det h)\chi(\det D)^{-1}\), equal by
\(\det h\equiv\det A\det D\bmod\pi^c\). Iwasawa decomposition
extends equality using their common parabolic law. The common value
at the identity is one. The same calculation in the smaller ranks
gives a nonzero common lower quotient vector.

Unprimed evaluation integrates the \(a\) entries of one nonpivot
column, setting the other removed nonpivot column to zero. Thus
\(e_{0,v}\phi=Q^{-ac}\phi_1\). For the primed calculation write
\(l=(s,t)\in K_0^\times\times K_0^\times\), the first factor
being the chosen place \(w\). Order the pivot rows and columns as
first, middle, last in the prescribed bases. In the conjugate row
polarization the position and modulation parts of
\(\bar l\,y\otimes x'\) are respectively the first-row last-pivot
entry \(s\) and the last-row first-pivot entry \(-t\).
Indeed switching from the ordinary to the conjugate row polarization
interchanges momentum and position with one minus sign; the source
pair has Hermitian off-diagonal entries \(1,1\), the target pair
has skew entries \(1,-1\), and conjugating \(l\) swaps its two
components. These coordinates occupy distinct matrix entries, so
the half-pairing in the Heisenberg action is zero.
\src{src/Development/WeilBound/ThetaLifts/LocalTheory/StandardValuations/AdelicInclusionMapPlaceSection.lean}{228}{The two entries for the specified Fourier signs}.

In this coordinate ordering the primed mixed evaluation integrates
the last row and the last pivot column in the middle rows. The
other removed entries vanish except for the position entry \(s\).
This is the matrix calculation of Lemma~\ref{tpd:mellin}, expressed
in the prescribed conjugate bases. The resulting pivot matrix is
\[
     \begin{pmatrix}0&0&s\\0&D_1&*\\ x&*&*\end{pmatrix},
       \qquad \det=-sx\det D_1 ,
\]
and the Fourier factor is \(\psi_v(-tx)\). The \(a\) integrated
nonpivot entries of the last row have joint volume \(Q^{-ac}\).
The other integrated entries, apart from \(x\), range independently
over \(\mathcal O\) of mass one. Since \(s,x\) and the entries
of \(D_1\) must be integral, the determinant is a unit precisely
when each of \(s,x,\det D_1\) is a unit. The minus sign in the
determinant is removed together with the Fourier minus sign by
\(x\mapsto-x\). The evaluated test at \(l\) is therefore
\[
 Q^{-ac}D_\chi(s)\widehat D_\chi(t)\phi'_1,\qquad
 \widehat D_\chi(t)=
       \int_{\mathcal O^\times}\chi(x)^{-1}\psi_v(tx)\,dx .
\]
The complementary place has character \(\chi^{-1}\), so the Mellin
weight is \(\chi(s)\chi(t)^{-1}|st|^{k+1}\).
Integration in \(s\) gives one. The Fourier transform vanishes
outside \(v(t)=-c\): below this shell additive integration
cancels on cosets modulo \(\pi^c\), and above it multiplicative
translation by a unit trivial on the additive factor but
nontrivial under \(\chi\) forces zero. At \(t=\pi^{-c}u\) its
value is \(\chi(u)Q^{-c}G(\chi,\psi_v)\) by finite residue
summation. Integration on that multiplicative shell therefore
gives \(\chi(\pi)^c Q^{ck}G(\chi,\psi_v)\).
The source compact averaging has no further effect on the chosen
test. The factors \(Q^{-ac}\) on the two sides cancel, and
Proposition~\ref{threal:localcone} extends the resulting scalar
identity to all matched upper tests. The nonvanishing of the
Gauss sum is Lemma~\ref{newtheta:gauss}.
\end{proof}

In the other descent Corollary~\ref{threal:ramA} gives the
corresponding factor \(\epsilon_v(1-k)\) with the same unit sign.

In this unprimed descent the integral line of \(l x\otimes y'\)
has coefficient ideal
\(\mathfrak b\,\mathfrak D_E^{-1}\bar{\mathfrak c}^{-1}\),
where \(\mathfrak D_E^{-1}\) is the absolute codifferent. Indeed
the trace-perfect target lattice has this ideal as its duality
range for the skew-Hermitian form; pairing \(y'\) against
\(\mathfrak c x'\) gives the indicated factor. Since \(E/F\)
is unramified at finite places the absolute codifferent extends
from \(F\), and its character value is one by
\(\nu|_{\mathbb A_F^\times}=1\). The two ideal exponents in the
unprimed branch are therefore \(+1,+1\), and the preceding
Fourier computation gives \(s=+1\) also here.

\begin{proposition}[The comparison at all relatively split good places]
\label{tpd:ordinary}
Let \(v\notin S\) split in \(E/F\), and choose either factor as
the first copy of \(K_0=F_v\). Use the integral-vacuum-normalized
local evaluations, compact mass one, and unit mass one on each
multiplicative group. Write \(z=k+1\) in the primed descent and
\(z=k\) in the unprimed descent. The lower quotient of the local
cone operator is the following scalar times the lower quotient of
the opposite zero evaluation, on every upper matched pair:
\[
   \nu_v(\mathfrak a_v)\,|\mathfrak a_v|_v^z L_{E,v}(\nu,z).
\]
Here \(\mathfrak a_v\) is the corresponding translation-line ideal
computed above, its modulus is that of an idele generating its two
local factors, and \(\nu_v\) is evaluated on such a generator.
The assertion includes \(\nu_w(\pi)=1\).
\end{proposition}
\begin{proof}
The source compact fixes the vacuum at this unramified split place:
in integral row bases its action is linear substitution with
unramified determinant character. In the two-plane evaluation the
translated vacuum maps to the smaller vacuum for
\(l\in\mathfrak a_v\) and to zero otherwise. To see the normalization,
use the two idempotent position and momentum groups, as in
Proposition~\ref{threal:ordinarysplit}. The vacuum is proportional
to the position lattice indicator, and its modulation lattice is
the character dual. Adaptedness decomposes both lattices into the
paired hyperbolic directions and the smaller tensor. Each exchanged
position entry is integrated by partial Fourier transformation at
zero. Translation in the nonintegrated entry leaves its indicator
nonzero precisely on that entry's lattice. Modulation on the
integrated entry gives its original integral on the character-dual
lattice and zero elsewhere. The two entries are distinct, so the
half-pairing vanishes also in residue characteristic two. Together
the two conditions say exactly that the displacement belongs to
the tensor lattice, equivalently \(l\in\mathfrak a_v\). All other
entries give the untranslated evaluation, normalized to the smaller
vacuum. This uses the character dual for the actual additive character,
including its codifferent. Integrating the two multiplicative shells
gives the displayed scalar. The opposite one-plane evaluation sends
vacuum to vacuum.

Let \(u=\nu_w(\pi)\) vary on the unit circle, keeping \(\beta_w\)
fixed and setting \(\alpha_w=\beta_w\nu_w\) with \(\nu_w\) unramified.
On the other split factor take the inverse characters. This family
is local. For \(u\) of infinite order the full source coinvariant
calculations of Propositions~\ref{split:kernel} and~\ref{split:good}
apply over \(K_0\) itself. Their proofs use the local orbital
functions and the infinite-order condition. The quotient-integral
and Jacquet argument of Proposition~\ref{threal:localcone} also uses
only those local facts and the split line decomposition. The smaller
zero evaluation is onto. Thus the two maps from the upper common
quotient to the smaller quotient are proportional for each such
\(u\). The vacuum calculation determines the scalar.

To extend the identity to \(u_0\) on the circle, realize the upper
and lower normalized induced targets by restriction of their functions
to the integral compact. The compact pictures are independent of
\(u\), since all varying inducing characters are unramified. Use
the orbital maps normalized to send both vacua to the spherical
vector of value one. For each fixed test they are continuous in
\(u\) on the circle. For the larger row map this follows by clearing
the determinant denominator as in Proposition~\ref{threal:ordinarysplit}.
The pivot-chart proof there works with the unramified value as a
variable: integration on a ball at zero gives a rational function
with denominator the indicated subproduct of
\(\prod_{j=0}^{n-1}(1-u Q^{j-d/2})\); a pivot of fixed valuation
contributes a Laurent monomial, and the other charts give the
lower-dimensional subproducts. Here \(Q\) is the residue size
and \(n\) the upper row rank. Clearing this denominator and normalizing
the lattice value gives Laurent polynomials in \(u\) up to a nonzero
fixed constant. A basis change with nonunit determinant introduces
a further Laurent monomial. The smaller row map is absolutely
convergent at its parameter and its spherical value has no zero or
pole on the circle. Use the corresponding dimensions and ranks in
the two lower maps. Compact-open invariance reduces their restrictions
to finitely many compact-picture evaluations at fixed level. Thus
continuity includes the normalized residue at a pole.

The zero evaluation is independent of \(u\). The local cone operator
also varies continuously on each test in a fixed finite-dimensional
span of smaller tests. The compact source action is independent of an
unramified change of character, and the restrictions to the two
displacement coordinates are compactly supported locally constant
functions with finitely many slices. On each compact open slice
away from zero the shell integrals are finite sums; near zero the
geometric series converge uniformly since \(z>0\). This applies
independently to the two local factors of \(E\).

For tests \(x,y\) with equal normalized upper images at \(u_0\)
denote the two upper maps by \(P_u,P'_u\). Choose a compact open
subgroup \(J\subset K_d\) fixing their images for every \(u\);
the upper compact actions on the tests are independent of \(u\).
The \(J\)-fixed space of the compact picture is finite-dimensional
and independent of \(u\). By surjectivity of the smaller-row map
\(P'_{u_0}\) from Proposition~\ref{threal:ordinarysplit}, choose
test preimages of a basis and average them over \(J\).
Surjectivity at the critical parameter uses the common irreducible
normalized unitary induction and the nonzero spherical image.
The matrix of \(P'_u\) on the span of these averaged preimages
stays invertible near \(u_0\); denote its inverse by \(H_u\).
The tests
\[
             y(u)=y+H_u(P_u x-P'_u y)
\]
are matched to \(x\) there, and converge to \(y\) in a fixed
finite-dimensional test space. For \(v=v_1\) the argument can
first be made on the full spaces and then restricted to supported
pairs. Apply the equality at infinite-order parameters tending
to \(u_0\) and pass to the limit on the two lower compact pictures.
The computed vacuum scalar is continuous as well. This only uses
unitary parameters.

Finally any allowed irreducible common quotient at this place is
isomorphic to the canonical one, and each map to it for trivial
source is unique up to scalar by the split local Howe theorem.
\cite[Theorem~1 and Corollary~6.3]{minguezHowe}
Equality of the two nonzero vacuum images forces the two scalar
changes to agree, in both the upper and lower models. Thus the
identity holds for every allowed choice of quotients.
\end{proof}

For comparison with the normalized recursion, denote the smaller
periods by \(I_1,I'_1\) and the primed, nonzero-cone test by
\(T\Phi'\). The comparison needed at the identity of the line
torus has three constants. With fixed models and measures they are
specified by the following identities, in which \(b,b'\) denote
the upper components with character \(\beta|\cdot|_E^{a/2}\):
\[
 b=C_0 I_1(e_0\Phi),\qquad b'=C_1 I'_1(T\Phi').
\]
The third requirement is that the following pair be matched on the
smaller spaces, with its second entry supported:
\[
 \big(C_2 F_X(-k)L^{S_0}(\nu,k+1)e_0\Phi,\ T\Phi'\big).
\]
If these identities hold and the smaller scalar is \(c_1\), then
\[
      C_0 c_1 b'=C_2 C_1 F_X(-k)L^{S_0}(\nu,k+1)b .
\]
Proposition~\ref{threal:endB} gives detection of the functions.
The upper scalar under the third requirement is therefore
\[
 c=\frac{c_1}{C F_X(-k)L^{S_0}(\nu,k+1)},\qquad
                  C=\frac{C_2C_1}{C_0}.
\]
Proposition~\ref{tpd:unfold} supplies the second identity for every
\(q>0\). For \(p>0\) on the unprimed side,
Proposition~\ref{threal:coneA} places the entire nonzero orbit in
the other character; the zero orbit gives the first identity by
evaluation, with a convergent smaller unprimed period by the
estimates of Proposition~\ref{threal:decay} with \(d\) decreased
by two. Thus \(C_0=1\) throughout for coherent mixed models. The
remaining test comparison is the third identity on the two smaller
matching quotients with one common normalization.
Propositions~\ref{tpd:ram} and~\ref{tpd:ordinary} give it at the
ramified and all relatively split unramified places, including the
unit sign \(s=+1\) and the factors of the two adapted line ideals.
The vacuum calculation and the partial Fourier transform supply the input for Propositions~\ref{tcr:descend} and \ref{tcr:unique}, which extend the inert identity to all smooth tests and compare the real operators on the polynomial module. Theorem~\ref{tcr:global} assembles their normalized tensor product.
\src{src/Development/WeilBound/ThetaLifts/LocalTheory/StandardValuations/ExistsVol.lean}{162}{Evaluation on a ramified matching pair};
\src{src/Development/WeilBound/ThetaLifts/LocalTheory/LocBOf.lean}{17}{Passage from these tests to local matching}.
The required \(C\) is common to the admissible characters and
the choices of finite matching maps.
\src{src/Development/AbelianVarietyTheory/Definitions/Divisors/CTSmallCone.lean}{239}{The two constant-term evaluations};
\src{src/Development/WeilBound/ThetaLifts/ConeHypotheses/KeyAOf.lean}{37}{Assembly of the three constants};
\src{src/Development/WeilBound/ThetaLifts/ConeHypotheses/PolDense.lean}{46}{The supported smaller matching identity}.
